\section{Discussion}
\label{sec:discussion}

\subsection{The Peat-Drainage Nexus and Hydrological Vulnerability}
\label{subsec:disc_nexus}

Our spatial and TreeSHAP findings confirm the link between artificial drainage and fire severity in degraded tropical peatlands \citep{Hooijer2012Biogeosciences, Ritzema2014HESS}. In pristine equatorial peat swamp forests, high water tables (typically within 10--20\,cm of the surface) and dense forest canopies maintain water saturation throughout the acrotelm, rendering peat domes virtually immune to deep fire ignition even during seasonal dry periods \citep{Wosten2008Geoderma, Dadap2021NatComm}. 

The excavation of extensive primary and secondary drainage canals across the Mega Rice Project (Blocks A, B, and C) permanently altered this hydrological buffer \citep{Silvius2005Wetlands, Ritzema2014HESS}. By piercing the elevated crests of ombrotrophic peat domes, canals induce continuous gravitational drawdown of phreatic water tables toward regional river systems. Our data reveal that 74.3\% of all active fire hotspots in Central Kalimantan ignited within delineated peatlands, with over 52.5\% situated in deep ($>2$\,m) and very deep ($>3$\,m) peat domes. Within the ex-PLG landscape, metric Euclidean distance to primary canals was the second most influential predictor of severe Fire Radiative Power ($\text{mean } |\text{SHAP}| = 0.0322$).

Canal drainage induces severe physical compaction, irreversible desiccation, and oxidative loss of the fibrous peat matrix \citep{Hooijer2012Biogeosciences, Graham2022SciRep}. When the water table drops below $-40$\,cm, aerobic microbial decomposition accelerates, transforming the uppermost meter of peat into a dry, porous fuel bed with exceptionally low bulk density and high air-filled porosity \citep{Usup2004Tropics}. Once ignited, combustion shifts from brief above-ground surface flaming to prolonged, oxygen-limited subterranean smoldering throughout the catotelm \citep{Wooster2018NatGeosci, Vetrita2020RemoteSensing}. The concentration of high-energy fires near canals indicates that drained corridors are primary ignition zones, requiring prioritized canal blocking and hydrological restoration.

\subsection{Resolving the Land Tenure Governance Void in Degraded Peatlands}
\label{subsec:disc_tenure}

Quantifying land tenure attribution shows that 75.42\% of all verified active fires (33,138 detections) occurred on unallocated Non-Concession land (\textit{Area Penggunaan Lain}, APL), while certified oil palm concessions (HGU) and commercial forestry licenses (PBPH) accounted for 16.56\% and 8.02\%, respectively.

These findings contrast with accounts that attribute Indonesian peat fires primarily to corporate concessions \citep{Tacconi2003CIFOR, Carlson2012PNAS}. At the same time, 24.58\% of fires (10,799 hotspots) occurred within legal corporate boundaries. Under Indonesian environmental jurisprudence (Law No.~32/2009 on Environmental Protection and Management, Article 88), concession holders bear strict liability (\textit{strict liability} / no-fault liability) for any fires occurring within their demarcated concessions, regardless of fault. The spatial identification of specific entities experiencing repeated fire clusters (e.g., PT.~Sangkowong Sinta with 954 hotspots; PT.~Rimba Makmur Utama with 964 hotspots) provides transparent, auditable evidence for corporate Environmental, Social, and Governance (ESG) compliance monitoring and legal enforcement.

However, the concentration of fires in unallocated APL zones points to an administrative gap across Central Kalimantan's degraded peatlands \citep{Cattau2016GCB, Putranto2023LandUsePol}. Following the formal abandonment of the Mega Rice Project in 1999, hundreds of thousands of hectares of drained peatland were designated as state land without assigned institutional managers or active concessionaires. In the absence of state enforcement, these lands have evolved into de facto open-access regimes characterized by speculative land acquisition, opportunistic clearance along canal banks, and accidental ignitions by informal resource extractors \citep{Cattau2016GCB}. Because no corporate entity is legally liable for fire suppression across unallocated APL lands, fire-fighting responsibilities fall on under-resourced regional disaster management agencies, forest fire suppression brigades, and volunteer community fire-fighting groups. Without formal management of these unallocated APL peatlands, regulating corporate concessions alone will not stop regional fire crises.

\subsection{Methodological Imperatives in Spatial Wildfire Machine Learning}
\label{subsec:disc_spatial_ml}

The comparison between Model~A and Model~B illustrates how spatial autocorrelation affects tree-based models in remote sensing. In published environmental machine learning literature, the inclusion of geodetic coordinates (\texttt{latitude} and \texttt{longitude}) as independent predictors is widely practiced under the presumption that spatial coordinates capture unobserved geographical heterogeneity \citep{Meyer2018EcolModel, Dormann2007Ecography}.

In Model~A, spatial coordinates acted as algorithmic shortcuts: decision trees recursively split coordinate axes to isolate historical spatial clusters. When evaluated using standard random train-test splitting (75:25), Model~A achieved an apparent ROC-AUC of 0.7956 because testing samples occupied the same spatial autocorrelation neighborhoods as training samples \citep{Tobler1970EconGeogr, Legendre1993Ecology}. However, when evaluated with Spatial Block Cross-Validation (\texttt{GroupKFold} by regency), Model~A's predictive capacity collapsed: mean ROC-AUC dropped to 0.5399, and the $F_1$-score fell to 0.0583, reaching 0.0000 across four of five spatial folds. The model was incapable of predicting severe fires in novel geographical regions because its learned decision boundaries relied on coordinate partitions that did not exist in out-of-sample blocks \citep{Roberts2017Ecography, Wadoux2021EcolModel}.

Conversely, by completely excluding geographic coordinates, Model~B forced the Random Forest ensemble to learn functional, physically grounded dependencies on antecedent precipitation deficits ($P_{\text{7d}}$), Euclidean canal proximity ($d_{\text{canal}}$), wind speed, and peat depth. While its random holdout ROC-AUC was slightly lower (0.7527), Model~B maintained positive $F_1$-scores across all five spatial blocks, achieving a mean spatial $F_1$-score of 0.1399, representing a 2.40-fold improvement over Model~A. These findings demonstrate that machine learning models in environmental remote sensing must be validated via spatial block cross-validation and constrained to their valid Area of Applicability \citep{MeyerPebesma2021NatComm, Schratz2019EcolModel}.

\subsection{Actionable Spatial Planning, Canal Zoning, and Transboundary Haze Dispersion}
\label{subsec:disc_spatial_planning_haze}

The empirical findings from this study directly inform applied spatial planning, hydrological zoning, and transboundary environmental risk governance across equatorial Southeast Asia.

First, the sharp non-linear elevation of fire risk within 3,000\,m of primary drainage channels identifies an empirical zoning threshold for ecological engineering. Provincial spatial planning (Rencana Tata Ruang Wilayah Provinsi / RTRWP) and the Peatland Ecosystem Protection and Management Plan (RPPEG) currently lack spatially differentiated hydrological buffer mandates for abandoned canal backbones. We propose that the Peatland and Mangrove Restoration Agency (BRGM) and the Ministry of Public Works prioritize a 3,000-m hydrological restoration corridor along the 13 primary canal axes in Mega Rice Project Blocks A and B. Within this active drainage ribbon, canal-blocking investments---incorporating cascading box dams and compacted mineral-peat dams---must be concentrated to maintain phreatic water tables above the regulatory $-40$\,cm threshold (PP No.~71/2014 and PP No.~57/2016), eliminating the primary ignition conduit for severe subterranean smoldering.

Second, resolving the spatial governance vacuum in unallocated Area Penggunaan Lain (APL) land (accounting for 75.42\% of active fires) requires transitioning abandoned state peatlands into legally administered tenurial regimes. Because APL peatlands lack institutional concessionaires or Forest Management Units (KPH), they function as open-access commons vulnerable to speculative burning. Transforming these territories into community-based Social Forestry schemes (\textit{Perhutanan Sosial}), such as Community Forests (Hutan Kemasyarakatan / HKm) or Village Forests (Hutan Desa / HD), provides local communities with formal land tenure rights tied to fire stewardship obligations, water table monitoring, and canal maintenance. Concurrently, corporate concessions (accounting for 24.58\% of ignitions) must be subject to automated spatial auditing dashboards that intersect satellite thermal anomalies with audited cadastral boundaries, enforcing strict liability under Law No.~32/2009 (Article 88) and incentivizing corporate fire prevention infrastructure.

Third, our analysis connects local peat combustion mechanics to regional transboundary air quality. During the dry season (August--October), prevailing Southwest Monsoon (Monsun Barat Daya) winds advect atmospheric emission plumes from southern Central Kalimantan north-northwestward across the Karimata Strait and South China Sea toward western Borneo (Kuching, Sarawak), the Singapore Strait, and Peninsular Malaysia (Kuala Lumpur) \citep{Gaveau2014PLOSONE, Kiely2021ACP}. In the October 2026 crisis, peak local surface fine particulate matter ($\text{PM}_{2.5}$) concentrations in Central Kalimantan reached $383.7\,\mu\text{g}/\text{m}^3$, with regional atmospheric dispersion simulations indicating downwind advection generating peak concentrations of $197.3\,\mu\text{g}/\text{m}^3$ in Kuching, $165.3\,\mu\text{g}/\text{m}^3$ in Singapore, and $191.7\,\mu\text{g}/\text{m}^3$ in Kuala Lumpur. These concentrations substantially exceed the World Health Organization (WHO) 24-hour guideline of $15\,\mu\text{g}/\text{m}^3$. By integrating biophysically grounded fire intensity predictions with meteorological monsoon trajectories, applied spatial geography provides the analytical foundation for anticipatory cross-border early warning systems under the ASEAN Agreement on Transboundary Haze Pollution.

\subsection{Study Limitations, Uncertainties, and Threats to Validity}
\label{subsec:disc_limitations}

Several methodological limitations and data uncertainties affect these findings:

\subsubsection{Sensor Resolution, Orbital Sampling Gaps, and Smoldering Combustion Physics}
First, active fire observations from the NASA FIRMS constellation (VIIRS 375\,m and MODIS 1\,km) are subject to satellite remote sensing constraints \citep{Schroeder2021RSE, Giglio2020RSE}. Both sensor systems operate on sun-synchronous, polar-orbiting platforms that pass over Central Kalimantan two to four times per day, with daytime overpasses concentrated between 12:30 and 14:30 local solar time. Consequently, diurnal fire front variations outside orbital passes, as well as nocturnal fire spread, escape observation. Persistent cloud cover and dense convective smoke plumes can also attenuate infrared radiance, introducing potential omission errors during peak burning periods \citep{Giglio2020RSE}.

In addition, infrared active fire algorithms detect radiance emitted by high-temperature surface flaming combustion (typically $>800$\,K). In degraded tropical peatlands, however, burning frequently occurs as subterranean smoldering combustion within the catotelm layer at temperatures of 500--700\,K without open flames \citep{Usup2004Tropics, Wooster2018NatGeosci, Vetrita2020RemoteSensing}. When smoldering burns beneath an intact forest canopy or deep within the peat matrix, emitted thermal energy is partially absorbed, reducing surface Fire Radiative Power. Therefore, the target threshold ($\text{FRP} \ge 12.0$\,MW) measures surface radiative power rather than the depth or mass of consumed peat.

\subsubsection{Scale Mismatch in Atmospheric Reanalysis and Potential Ecological Fallacy}
Second, meteorological forcing derived from ECMWF ERA5-Land reanalysis features a nominal spatial resolution of $0.1^\circ \approx 9$\,km, aggregated at the administrative regency level on a daily cadence. In extensive jurisdictions such as Kapuas Regency (spanning over 15,000\,km$^2$ from northern uplands to the southern coast), assigning uniform daily meteorological values introduces an unavoidable scale mismatch. Localized convective thunderstorms that deposit heavy precipitation over an isolated peat dome cannot be resolved at this spatial scale. Future operational implementations should integrate dynamically downscaled weather radar data or dense automated weather station (AWS) networks.

\subsubsection{Temporal Scope of the Acute Fire Crisis}
Third, the machine learning models were trained and cross-validated on an eight-day fire crisis window (5--12 October 2026, comprising 43,937 observations) characterizing the transition from severe drought to monsoonal rainfall. While this sample captures extreme fire behavior under dry fuel conditions, the empirical parameter weights reflect this specific climatological pulse. Model application across the early dry season (June--July) or during multi-year wet regimes (e.g., strong La Ni\~na cycles) requires dynamic seasonal retraining and domain adaptation.

\subsubsection{Spatial Transferability, Covariate Shift, and Class Imbalance Trade-offs}
Fourth, while Model~B prevented spatial memorization collapse, its mean Spatial Block CV ROC-AUC averaged 0.4950 ($\pm$0.0490) and its spatial CV $F_1$-score reached 0.1399. This reduction reflects spatial covariate shift between disparate administrative blocks. For example, primary canal drainage networks are concentrated in Pulang Pisau and Kapuas (ex-PLG Blocks A and B), whereas Seruyan Regency features distinct, uncanalized coastal peat domes with different hydrological conditions. The class imbalance (15.13\% severe fire class; 1:5.6 ratio) combined with inverse class weighting (\texttt{class\_weight=balanced}, $w_1 = 3.30$) prioritized recall over precision. On the independent holdout test set ($n = 10,985$), Model~B captured 65.82\% of all severe fires ($\text{TP} = 1,094$), with 2,737 false positives ($\text{FP}$), yielding an extreme-class precision of 28.56\% and an overall holdout $F_1$-score of 0.3983.

In wildfire disaster management, this trade-off reflects loss asymmetry: failing to detect a severe fire before it penetrates the subterranean peat column incurs severe ecological and economic losses, whereas a false alarm triggers precautionary reconnaissance. However, excessive false alarms risk alert fatigue among emergency response crews. To operationalize the model safely, we recommend a three-tier risk stratification protocol based on predicted probabilities:
\begin{itemize}
    \item Advisory Tier ($p \in [0.30, 0.50)$): Heightened community patrolling and water table monitoring.
    \item Watch Tier ($p \in [0.50, 0.70)$): Ground dispatch of forest fire response teams to test surface peat moisture.
    \item Warning Tier ($p \ge 0.70$): Immediate deployment of aerial water-bombing aircraft and activation of canal-blocking control gates.
\end{itemize}

\subsubsection{Cadastral Boundary Resolution and APL Land Tenure Heterogeneity}
Finally, the attribution of 75.42\% of active fires to unallocated Non-Concession (APL) land represents an aggregate legal classification. In practice, APL land in Central Kalimantan encompasses a diverse mosaic: customary ancestral claims, transmigration farming settlements, independent smallholder oil palm plots with formal freehold titles, and abandoned, degraded state peatlands \citep{Putranto2023LandUsePol}. Disentangling these sub-categories requires fine-scale participatory cadastral mapping and village-level tenure boundary demarcation.
